\documentclass[10pt,twocolumn]{article}

\usepackage[margin=0.75in]{geometry}
\usepackage{times}
\usepackage{microtype}
\usepackage{amsmath,amssymb,mathtools}
\usepackage{booktabs}
\usepackage[numbers,sort&compress]{natbib}
\usepackage[colorlinks=true,citecolor=blue,linkcolor=blue,urlcolor=blue]{hyperref}

\title{Related Work for Cone-Aware Primal-Dual Graph Unrolling}
\author{}
\date{}

\begin{document}
\maketitle

\begin{abstract}
This note summarizes related work on sparse optimization, graph neural networks for mathematical programming, and algorithm-unrolled learning-to-optimize methods.
The literature contains mature unrolled solvers for LASSO-type sparse recovery, WL-based GNN expressivity results for LP/QP/QCQP/SOCP, and PDHG-unrolled architectures for LPs and QPs.
The remaining gap is a cone-aware, dimension-agnostic primal-dual graph architecture that unrolls SOCP first-order iterations while retaining explicit second-order cone projections.
\end{abstract}

\section{Related Work}

\paragraph{GNN expressivity for learning to optimize.}
Graph neural networks have become a common inductive bias for mathematical programs whose variables and constraints admit permutation-equivariant representations.
The WL-based line of work represents an optimization instance as a graph and relates the expressive power of message passing to Weisfeiler--Lehman graph tests.
For LPs, \citet{chen2022lpgnn} represent variables, constraints, and coefficients by a weighted bipartite graph and establish approximation results for LP feasibility, boundedness, optimal values, and optimal solutions.
This framework has been extended to quadratic and mixed-integer quadratic programs \citep{chen2024qpgnn}, as well as convex quadratically constrained quadratic programs \citep{wu2024qcqpgnn}.
More general provable GNN solvers for convex optimization have also been studied \citep{qian2025provableconvexgnn}.
Recent SOCP-GNN results extend WL-based expressivity to second-order cone constraints by augmenting the graph representation to encode conic blocks and prove universality and sample-complexity guarantees \citep{li2025socpgnnworkshop,li2026socpgnn}.
These results show that SOCPs can be represented by suitably designed GNNs, but they do not by themselves specify an algorithmic trajectory, a residual-driven primal-dual loss, or an explicit solver warm-start procedure.
Complementary negative results for SDPs further indicate that graph architectures must match the algebraic structure of the underlying cone; standard message passing may fail when the cone geometry introduces higher-order symmetries \citep{qian2026sdpgnn}.

\paragraph{Algorithm-unrolled learning to optimize.}
Algorithm unrolling maps iterations of an optimization algorithm to neural network layers, preserving useful solver structure while replacing hand-designed parameters by trainable ones \citep{monga2021algorithm}.
The primal-dual hybrid gradient method \citep{chambolle2011primaldual} is especially relevant because it solves saddle-point formulations of the form
\begin{equation}
\begin{aligned}
    \min_x \; & f(x)+g(Kx) \\
    \Longleftrightarrow\quad
    \min_x\max_y \; & f(x)+\langle Kx,y\rangle-g^*(y),
\end{aligned}
    \label{eq:pdhg_saddle}
\end{equation}
PDHG-Net unrolls this structure for large-scale LPs and uses channel expansion to avoid dense trainable matrices tied to a fixed problem size \citep{li2024pdhgnet}.
PDQP-Net adopts a related primal-dual unrolling strategy for convex QPs and trains with KKT-informed unsupervised losses \citep{yang2024pdqpnet}.
Other constrained unrolling approaches use coupled primal and dual GNNs to emulate dual-ascent dynamics for general constrained problems \citep{hadou2025unrolledgnn}, while doubly stochastic PDHG studies blockwise primal-dual updates and sampling mechanisms in classical first-order optimization \citep{xiao2026dspdhg}.
This AU-based paradigm provides algorithmic interpretability and size generalization, but extending it from LP/QP updates to conic programs requires representing nonlinear cone projections and cone-block interactions inside the learned architecture.

\paragraph{Sparse optimization and graph unrolling.}
Sparse recovery provides a canonical example of successful unrolling.
The LASSO \citep{tibshirani1996lasso} solves
\begin{equation}
    \min_{x\in\mathbb{R}^n}
    \frac{1}{2}\|Ax-b\|_2^2 + \lambda\|x\|_1 ,
    \label{eq:lasso}
\end{equation}
and proximal-gradient methods such as ISTA and FISTA exploit the soft-thresholding proximal operator of the $\ell_1$ norm \citep{beck2009fista}.
LISTA unrolls ISTA into a trainable architecture for sparse coding \citep{gregor2010lista}, and subsequent work established convergence guarantees and efficient hyperparameterized variants \citep{chen2018linearlista,chen2021hyperlista}.
Graph unrolling networks extend the same idea to graph signal denoising and include a graph sparse-coding architecture, whose objective can be written as a LASSO-type problem with a graph-filter dictionary \citep{chen2021graphunrolling}.
Thus, sparse optimization already has mature neural and graph-unrolled variants.
The unresolved LASSO setting considered here is not sparse coding on a fixed graph dictionary, but graph-parametric optimization over arbitrary sensing matrices, where the matrix $A$ itself defines a variable-measurement bipartite graph and the learned update rule must generalize across dimensions.

\paragraph{Conic optimization and SOCP first-order methods.}
SOCPs extend LPs and convex QPs by introducing constraints of the form
\begin{equation}
\begin{aligned}
    \min_x \quad & c^\top x \\
    \text{s.t.}\quad & Fx \leq g,\quad l\leq x\leq r,\\
    & \|A_i x+b_i\|_2 \leq d_i^\top x+e_i,\quad i=1,\ldots,q .
\end{aligned}
\label{eq:socp}
\end{equation}
where $\|A_i x+b_i\|_2 \leq d_i^\top x+e_i$ is equivalent to membership in the second-order cone $\mathcal{K}_{\mathrm{soc}}=\{(t,z):\|z\|_2\leq t\}$ \citep{lobo1998socp}.
Classical primal-dual and splitting methods handle such constraints through projections onto products of cones.
Recent large-scale conic solvers build on restarted PDHG, adaptive step-size rules, GPU-accelerated sparse matrix-vector products, and efficient projections onto zero, nonnegative, second-order, rotated second-order, exponential, and semidefinite cones \citep{lin2026pdcs}.
These solvers establish that PDHG is a practical algorithmic backbone for conic programs.
However, they remain hand-designed first-order methods rather than learned, dimension-agnostic graph architectures.

\paragraph{Scope of the present work.}
The closest learning-based SOCP work is WL-based: it proves that specialized GNNs can represent SOCP feasibility and solutions \citep{li2026socpgnn}.
The closest unrolled learning-to-optimize works are AU-based: they unroll PDHG/PDQP for LPs and QPs using channel expansion and primal-dual residual losses \citep{li2024pdhgnet,yang2024pdqpnet}.
The closest classical solver work applies PDHG directly to conic programs with explicit cone projections \citep{lin2026pdcs}.
To our knowledge, existing work has not combined these three ingredients into a cone-aware, PDHG-unrolled GNN for SOCPs that retains explicit second-order cone projections while learning step sizes, extrapolation, block-sampling probabilities, or channel-mixing maps.

\begin{table*}[t]
\centering
\caption{Taxonomy of closely related learning-to-optimize paradigms.}
\label{tab:novelty}
\small
\begin{tabular}{@{}p{0.18\textwidth}p{0.36\textwidth}p{0.40\textwidth}@{}}
\toprule
Paradigm & Representative coverage & Structure not covered \\
\midrule
Sparse unrolling & LISTA, HyperLISTA, graph sparse-coding unrolling & Arbitrary matrix-defined LASSO graphs with size transfer \\
WL-based GNNs & LP, QP, QCQP, and SOCP expressivity & Algorithmic primal-dual trajectories and solver warm starts \\
AU-based L2O & PDHG-Net for LPs and PDQP-Net for QPs & Explicit SOC projection and cone-block message passing \\
Conic first-order solvers & PDHG-based matrix-free conic solvers & Learned graph/channel parameterization across instances \\
General constrained GNNs & Dual-ascent unrolled GNNs & SOCP-specific PDHG updates and cone geometry \\
\bottomrule
\end{tabular}
\end{table*}

\bibliographystyle{plainnat}
\bibliography{../bib/references}

\end{document}
